\begingroup\footnotesize
\begin{longtable}{@{}p{.22\linewidth}p{.74\linewidth}@{}}\toprule
\textbf{Término (cláusula)} & \textbf{Definición parafraseada y ejemplo aplicado a Spotify}\\\midrule\endhead
Requisito (3.1.19) & Enunciado que expresa una necesidad con sus condiciones y restricciones. Ejemplo: poder añadir una canción a una lista propia.\\
Condición (3.1.6) & Circunstancia que limita cuándo aplica el requisito. Ejemplo: búsqueda con conexión establecida.\\
Restricción (3.1.7) & Límite impuesto externamente a la solución. Ejemplo: descarga musical restringida por tipo de cuenta.\\
Contexto de uso (3.1.8) & Usuarios, tareas, medios y entorno de utilización. Ejemplo: estudiante que escucha música desde Windows.\\
Interesado (3.1.28) & Persona u organización con expectativas legítimas. Ejemplo: oyente suscrito a Premium.\\
Elicitación (3.1.20) & Técnicas para descubrir necesidades y registrarlas. Ejemplo: entrevistar usuarios sobre playlists.\\
Ingeniería de requisitos (3.1.21) & Disciplina que descubre, analiza, acuerda y mantiene requisitos. Ejemplo: formalizar necesidades de oyentes.\\
Trazabilidad de requisitos (3.1.23) & Registro del origen y derivación de los requisitos hacia otros niveles. Ejemplo: vincular una necesidad de organizar música con R-01 y sus requisitos derivados.\\
Validación de requisitos (3.1.25) & Comprobación de que los requisitos describen el sistema que los interesados realmente necesitan. Ejemplo: confirmar la utilidad de R-01.\\
Verificación de requisitos (3.1.26) & Examen de la redacción y organización de los requisitos para determinar si están bien formados. Ejemplo: detectar «muy rápido» en R-11.\\
\bottomrule\end{longtable}\endgroup
